% ==============================================================================
% T0 Theory: Document 162 (English)
% Title: T0 Optimization Strategies: QM-Optimization and Coherent Ising Machine Compared
% Author: Johann Pascher
% Date: February 2026
% ==============================================================================

\section*{Abstract}

The T0 theory contains two complementary optimization frameworks: Document~034
(\emph{QM-Optimization}) describes how quantum computers can be optimized through
geometric transformation in cylindrical phase space and through noise suppression
via time-field modulation. Document~161 (\emph{T0 Ising Machine}) shows that
combinatorial NP-hard optimization can be solved through geometric relaxation of the
mass field $\mfield$. This document~162 systematically compares both approaches and
demonstrates: they are not competing, but \emph{complementary levels of the same
geometric principle} -- optimization as physical relaxation into a geometric minimum.

\clearpage

\section{Introduction: Two Optimization Levels in T0}
\label{sec:intro}

\subsection{The Starting Problem}

Quantum computing struggles with noise and suboptimality on two fundamentally different
levels:

\begin{enumerate}
    \item \textbf{Gate level}: Individual quantum gates are disturbed by decoherence,
    phase noise, and fractal damping. The question: how does one optimize
    \emph{individual operations}?

    \item \textbf{Problem level}: Combinatorial optimization problems (NP-hard) are
    hard to solve even with perfect gates. The question: how does one find
    \emph{global minima} in vast solution spaces?
\end{enumerate}

The T0 theory addresses \textbf{both levels} -- but with different tools and at
different levels of abstraction. Document~034 \cite{pascher:qmopt} solves Problem~1;
Document~161 \cite{pascher:ising} solves Problem~2.

\subsection{The Common Root}

Both approaches share the same geometric origin: the parameter
$\xipar = 4/30000 \approx 1.333 \times 10^{-4}$
and the time-field binding condition $T \cdot m = 1$.
The differences are not conceptual -- they are differences in the
\emph{scaling level} of application.

\clearpage

\section{Document 034: QM-Optimization through Geometric Formalism}
\label{sec:doc034}

\subsection{The Cylindrical Phase Space}

In Document~034 \cite{pascher:qmopt}, Hilbert space is replaced by a
\textbf{cylindrical phase space}. A qubit is not a 2D complex vector, but a
point $(z, r, \theta)$:

\begin{itemize}
    \item $z \in [-1, 1]$: classical basis projection ($z=+1 \equiv |0\rangle$,
    $z=-1 \equiv |1\rangle$)
    \item $r \geq 0$: coherence radius (superposition magnitude), with $z^2 + r^2 = 1$
    \item $\theta \in [0, 2\pi)$: relative phase of the superposition
\end{itemize}

Gate operations are geometric transformations of this space, not matrix multiplications.

\subsection{Fractal Damping as Intrinsic Noise}

The X-gate (bit-flip) is not an ideal $\pi$-rotation in the T0 formalism, but a
damped rotation by angle $\alpha = \pi \cdot \Kfrak$:

\begin{align}
    z' &= z \cos(\alpha) - r \sin(\alpha) \\
    r' &= z \sin(\alpha) + r \cos(\alpha)
\end{align}

with the fractal damping constant $\Kfrak = 1 - 100\xipar$.
A perfect flip ($\alpha = \pi$) is physically impossible -- spacetime inherently
damps every rotation. \textbf{This is a testable prediction.}

\subsection{The Three Optimization Strategies from Doc. 034}

\begin{enumerate}
    \item \textbf{T0-Topology-Compiler}: Qubits are arranged not to minimize SWAP
    operations, but to minimize the cumulative \emph{fractal path length} of all
    entangling operations. Critically interacting qubits must be brought physically
    closer together.

    \item \textbf{Harmonic resonance frequencies}: Qubit frequencies are calibrated
    to ``magic'' frequencies arising from the harmonic structure of the T0 time field:
    \begin{equation}
        f_n = \left(\frac{\Ezero}{h}\right) \cdot \xipar^2 \cdot (\phiT^2)^{-n}
        \label{eq:harmonic_freq}
    \end{equation}
    For superconducting qubits, sweet spots emerge at approximately \textbf{6.24 GHz}
    ($n=14$) and \textbf{2.38 GHz} ($n=15$). This calibration intrinsically reduces
    phase noise.

    \item \textbf{Active time-field modulation}: A high-frequency
    \emph{time-field pump} irradiates idle qubits to average out the fundamental
    $\xipar$-noise and overcome the passive $T_2$ coherence limit -- active
    decoherence suppression instead of passive waiting.
\end{enumerate}

\clearpage

\section{Document 161: T0 Ising Machine -- Combinatorial Optimization}
\label{sec:doc161}

\subsection{The Ising Model as Optimization Language}

In Document~161 \cite{pascher:ising}, it is shown that NP-hard combinatorial
problems can be formulated as minimization of the Ising Hamiltonian:

\begin{equation}
    H_{\text{Ising}} = -\frac{1}{2} \sum_{i \neq j} J_{ij} \, \sigma_i \sigma_j
    - \sum_i h_i \, \sigma_i
    \label{eq:ising}
\end{equation}

A Coherent Ising Machine (CIM) finds the ground state of this Hamiltonian through
physical relaxation of laser pulses (spins = phase states 0° or 180°).

\subsection{The T0 Ising Machine: Emergent Couplings}

The T0 theory extends the CIM fundamentally: the couplings $J_{ij}$ are not
externally programmable -- they geometrically emerge from the time field:

\begin{equation}
    J_{ij}^{(\text{T0})} = \int \Tfield_i \cdot \Tfield_j \cdot G(x_i, x_j) \, d^3x
    \label{eq:t0_coupling}
\end{equation}

The energy functional is:
\begin{equation}
    H_{\text{Ising}} \leftrightarrow E_{\text{T0}} = \int |\nabla \Tfield|^2 \, d^4x
    \label{eq:energy_equiv}
\end{equation}

\subsection{Noise in the CIM: The Annealing Problem}

In the CIM, noise plays a paradoxical role: it is simultaneously problem and tool.
Too little thermal noise leads to local minimum traps; too much noise prevents
convergence. The annealing temperature parameter $T_{\text{ann}}$ balances this
conflict.

In the T0 Ising Machine, $\Kfrak^{3/2}$ takes this role -- not as an external
parameter, but as a geometrically determined fractal correction factor.

\clearpage

\section{Systematic Comparison: Doc. 034 vs. Doc. 161}
\label{sec:comparison}

\subsection{Overview Table}

\begin{table}[h]
    \centering
    \caption{Direct comparison of the two T0 optimization frameworks}
    \label{tab:main_comparison}
    \resizebox{\textwidth}{!}{\begin{tabular}{@{}llll@{}}
        \toprule
        Aspect & Doc.\ 034: QM-Optimization & Doc.\ 161: T0 Ising Machine & Common Root \\
        \midrule
        Target problem & Gate noise, decoherence & NP-hard optimization & Geometric relaxation \\
        Scaling level & Single qubit / gate & Full system ($N$ spins) & $\xipar = 4/30000$ \\
        Spin representation & $(z, r, \theta)$ in cylinder & Phase $0°/180°$ (CIM) & $T \cdot m = 1$ duality \\
        Role of noise & Disturbance (minimize) & Tool (annealing) & $\Kfrak$ correction \\
        Coupling & Gate pulses (external) & $J_{ij}$ (geometrically emergent) & Time field $\Tfield$ \\
        Optimization goal & Maximum gate fidelity & Global minimum of $H$ & Energy minimization \\
        Determinism & Explicit (geometric trafo) & Hidden (OPO equations) & Both deterministic \\
        Testable prediction & Damped X-flip & $\ln(N)$ residual after CIM & $\xipar$ signatures \\
        \bottomrule
    \end{tabular}}
\end{table}

\subsection{Parallel 1: Noise -- Enemy or Ally?}

This is the deepest conceptual difference between both documents:

\begin{table}[h]
    \centering
    \caption{The opposing role of noise in both frameworks}
    \label{tab:noise_role}
    \resizebox{\textwidth}{!}{\begin{tabular}{@{}lll@{}}
        \toprule
        Aspect & Doc.\ 034: QM-Optimization & Doc.\ 161: T0 Ising Machine \\
        \midrule
        Noise is... & \textbf{Enemy} -- disturbs gate fidelity & \textbf{Tool} -- enables annealing \\
        T0 cause & Fractal damping $\Kfrak$ & Thermal fluctuations in OPO \\
        Strategy & Actively suppress & Controlledly deploy \\
        Mechanism & Time-field pump modulates $T(x)$ & Temperature parameter $T_{\text{ann}}$ \\
        Optimum & Noise $\rightarrow 0$ & Noise $\rightarrow$ minimum (then $\rightarrow 0$) \\
        T0 analogue & $\Kfrak \rightarrow 1$ (perfect gates) & $\Kfrak^{3/2}$ as annealing force \\
        \bottomrule
    \end{tabular}}
\end{table}

\begin{important}[Resolution of the Apparent Contradiction]
    The apparent contradiction -- noise is bad in Doc.\ 034, good in Doc.\ 161 --
    dissolves through the scaling level: at the gate level, any noise is harmful.
    At the system level, controlled noise is necessary to escape local minima.
    T0 describes the same physical mechanism ($\Kfrak$) at both levels.
\end{important}

\subsection{Parallel 2: Fractal Damping as Universal T0 Mechanism}

\begin{table}[h]
    \centering
    \caption{$\Kfrak$ as a universal T0 mechanism at different levels}
    \label{tab:kfrak_universal}
    \resizebox{\textwidth}{!}{\begin{tabular}{@{}llll@{}}
        \toprule
        Context & Doc.\ 034 & Doc.\ 161 & Physical Meaning \\
        \midrule
        $\Kfrak$ definition & $1 - 100\xipar$ & $K_{\text{frak}}^{3/2}$ (g-2 correction) & Fractal spacetime structure \\
        Acts on... & X-gate rotation (damps) & Annealing dynamics (regulates) & Deviation from ideal \\
        Mathematically & $\alpha = \pi \cdot \Kfrak$ & $\mathcal{D}(N) = e^{-\xipar \ln N/\Dfrak}$ & Geometric damping \\
        Measurable via & Residual flip error & $\ln(N)$ residual after CIM & $\xipar$ signatures \\
        Testable on & IBM quantum computer & CIM with variable spin count & Both platforms \\
        \bottomrule
    \end{tabular}}
\end{table}

\subsection{Parallel 3: Frequency Optimization vs.\ Topology Optimization}

Both documents propose optimizations of the \emph{system architecture} -- at
different levels:

\begin{table}[h]
    \centering
    \caption{Architecture optimizations: gate level vs.\ system level}
    \label{tab:architecture}
    \resizebox{\textwidth}{!}{\begin{tabular}{@{}lll@{}}
        \toprule
        Optimization type & Doc.\ 034: QM-Optimization & Doc.\ 161: T0 Ising Machine \\
        \midrule
        Frequency optimization & Harmonic resonance frequencies $f_n$ & -- \\
        Topology optimization & T0-Topology-Compiler (fractal path length) & 4D torus topology \\
        Decoherence protection & Active time-field modulation & Toroidal boundary conditions \\
        Coupling structure & Gate pulses pre-compensated & $J_{ij}$ geometrically emergent \\
        Goal & Max.\ gate fidelity & Min.\ Ising Hamiltonian \\
        \bottomrule
    \end{tabular}}
\end{table}

\subsection{Parallel 4: Determinism and Measurability}

\begin{table}[h]
    \centering
    \caption{Determinism and testable predictions of both documents}
    \label{tab:determinism}
    \resizebox{\textwidth}{!}{\begin{tabular}{@{}lll@{}}
        \toprule
        Aspect & Doc.\ 034: QM-Optimization & Doc.\ 161: T0 Ising Machine \\
        \midrule
        Type of determinism & Explicit: geometric transformations & Hidden: OPO field equations \\
        Source of stochasticity & Fractal spacetime fluctuations & Quantum vacuum initial conditions \\
        Testable prediction 1 & Damped X-flip: $\alpha < \pi$ & $\ln(N)$-scaling residual \\
        Testable prediction 2 & Sweet spots at 6.24/2.38 GHz & $J_{ij}$ deviations from free param. \\
        Platform & IBM Brisbane/Sherbrooke & CIM with 1000+ spins \\
        Quantification & $\Delta\alpha = \pi(1 - \Kfrak) \approx 10^{-2}$ & $\xipar \approx 1.3 \times 10^{-4}$ \\
        \bottomrule
    \end{tabular}}
\end{table}

\clearpage

\section{New Optimization Opportunities through the Connection}
\label{sec:new_opportunities}

The systematic comparison of both documents opens optimization opportunities not
visible in either document alone:

\subsection{Hybrid Optimization: Gate Level + System Level}

\begin{keyresult}[Hybrid T0 Optimization Strategy]
    A fully T0-optimized quantum architecture combines both levels:
    \begin{enumerate}
        \item \textbf{Gate level} (Doc.\ 034): Harmonic frequency calibration +
        time-field pump for maximum gate fidelity
        \item \textbf{System level} (Doc.\ 161): T0 Ising Machine for combinatorial
        optimization with geometrically emergent couplings
        \item \textbf{Linking level}: The T0-Topology-Compiler (Doc.\ 034) minimizes
        fractal path lengths -- this is the same geometric principle as the torus
        structure in Doc.\ 161
    \end{enumerate}
\end{keyresult}

\subsection{Noise as a Bridge between Both Levels}

The dualistic role of noise opens a new optimization strategy:

\begin{itemize}
    \item \textbf{Phase 1 (Annealing)}: Controlled noise ($\Kfrak$ regime)
    -- the T0 Ising Machine searches for global minima
    \item \textbf{Phase 2 (Precision)}: Noise suppression (time-field pump)
    -- quantum gates execute the solution with maximum fidelity
    \item \textbf{Transition}: The fractal parameter $\Kfrak$ regulates both phases --
    as annealing mechanism (Doc.\ 161) and as gate damping (Doc.\ 034)
\end{itemize}

\begin{equation}
    \text{Optimization} = \underbrace{\Kfrak^{3/2}\text{-annealing}}_{\text{Doc.\ 161:\ global minimum}}
    + \underbrace{\Kfrak \rightarrow 1\text{-precision}}_{\text{Doc.\ 034:\ gate fidelity}}
    \label{eq:hybrid_optimization}
\end{equation}

\subsection{New Testable Prediction: Correlation between Gate Error and CIM Residual}

If $\Kfrak$ describes the same universal T0 parameter at both levels, then:

\begin{equation}
    \frac{\Delta\alpha_{\text{Gate}}}{\pi} = \frac{\mathcal{D}_{\text{CIM}}(N)}{\ln N}
    = \xipar + O(\xipar^2)
    \label{eq:universal_kfrak}
\end{equation}

This is a \textbf{new quantitative prediction}: the measured gate flip error on IBM
quantum computers should converge with the $\ln(N)$ residual in CIM measurements to
\emph{the same numerical value} $\xipar \approx 1.3 \times 10^{-4}$.
This is experimentally verifiable with current technology.

\clearpage

\section{Summary: Complementary Levels of the Same Principle}
\label{sec:summary}

\begin{table}[h]
    \centering
    \caption{Full synthesis: both documents as complementary T0 levels}
    \label{tab:synthesis}
    \resizebox{\textwidth}{!}{\begin{tabular}{@{}llll@{}}
        \toprule
        Dimension & Doc.\ 034 & Doc.\ 161 & Connection \\
        \midrule
        Level & Gate / qubit & System / spins & Same $\xipar$ \\
        Noise & Enemy (suppress) & Tool (control) & $\Kfrak$ at both levels \\
        Representation & Cylindrical $(z,r,\theta)$ & Phase $0°/180°$ & $T \cdot m = 1$ states \\
        Topology & Fractal path length & 4D torus & Geometric minimization \\
        Coupling & Gate pulses (pre-compensated) & $J_{ij}$ (emergent) & Time field $\Tfield$ \\
        Prediction & $\Delta\alpha \approx 10^{-2}$ & $\ln(N)$ residual & $\xipar = 4/30000$ \\
        \bottomrule
    \end{tabular}}
\end{table}

The central insight is:

\begin{revolutionary}[The Unified T0 Optimization Perspective]
    Both gate noise (Doc.\ 034) and combinatorial NP-optimization (Doc.\ 161) are
    manifestations of \emph{the same geometric principle}: physical relaxation into
    the minimum of the T0 energy functional
    $E_{\text{T0}} = \int |\nabla \Tfield|^2 \, d^4x$.
    The difference is not conceptual -- it is the scaling level of observation.
    The fractal correction $\Kfrak$ appears at both levels as the same universal
    parameter with numerically consistent predictions.
\end{revolutionary}

\begin{thebibliography}{9}

\bibitem{pascher:qmopt}
J.~Pascher,
\newblock T0 Quantum Mechanics: Geometric Formalism and Optimization Strategies
for Quantum Computers,
\newblock T0 Documentation Series, Document 034 (2025).

\bibitem{pascher:ising}
J.~Pascher,
\newblock T0 Theory and Coherent Ising Machines: Structural Parallels,
\newblock T0 Documentation Series, Document 161 (2026).

\bibitem{pascher:fundamentals}
J.~Pascher,
\newblock T0 Theory: Time-Mass Duality as a Fundamental Principle,
\newblock T0 Documentation Series, Document 003 (2025).

\bibitem{pascher:g2}
J.~Pascher,
\newblock Fractal Corrections and the g-2 Anomaly,
\newblock T0 Documentation Series, Document 018 (2025).

\bibitem{pascher:quantum}
J.~Pascher,
\newblock T0 Quantum Computing: Qubits and Logic Gates,
\newblock T0 Documentation Series, Document 147 (2025).

\bibitem{pascher:scramblons}
J.~Pascher,
\newblock T0 Projection onto Scramblon Physics,
\newblock T0 Documentation Series, Document 148 (2026).

\bibitem{inagaki_cim_2016}
T.~Inagaki et al.,
\newblock A coherent Ising machine for 2000-node optimization problems,
\newblock \emph{Science} \textbf{354}, 603 (2016).

\end{thebibliography}
